% ===== PRÉAMBULE CANONIQUE — à coller VERBATIM en tête de CHAQUE .tex =====
\documentclass[11pt,a4paper]{article}
\usepackage[utf8]{inputenc}\usepackage[T1]{fontenc}\usepackage{lmodern}
\usepackage[french]{babel}
\usepackage{amsmath,amssymb,mathtools}\usepackage{icomma}
\usepackage{siunitx}\sisetup{locale=FR}  % \qty{}{}, \si{}, \ang{}
\usepackage{xcolor}\usepackage{colortbl}
\usepackage{tikz}\usepackage{pgfplots}\pgfplotsset{compat=1.18}
\usepackage[siunitx,europeanresistors]{circuitikz} % circuits (élec.)
\usepackage[most]{tcolorbox}\usepackage{enumitem}\usepackage{array,booktabs}
\usepackage[a4paper,margin=2.2cm]{geometry}
\usetikzlibrary{arrows.meta,calc,angles,quotes,positioning,patterns,%
  backgrounds,shapes.geometric,decorations.pathmorphing,decorations.markings,intersections}
\definecolor{primary}{RGB}{222,49,99}\definecolor{primarydark}{RGB}{150,22,55}
\definecolor{defcol}{RGB}{0,114,178}\definecolor{methcol}{RGB}{52,92,148}
\definecolor{excol}{RGB}{0,135,90}\definecolor{attncol}{RGB}{200,40,55}
\tcbset{boxbase/.style={enhanced,breakable,boxrule=0.9pt,arc=2.5pt,
  left=3mm,right=3mm,top=2.3mm,bottom=2.3mm,fonttitle=\bfseries,coltitle=white}}
% --- boîtes TUTORIEL (noms EXACTS, ne pas renommer) ---
\newtcolorbox{cle}[1][]{boxbase,colback=defcol!6,colframe=defcol,
  title={\textsc{À retenir}\ifx\relax#1\relax\relax\else\textmd{ : \itshape #1}\fi}}
\newtcolorbox{methode}[1][]{boxbase,colback=methcol!7,colframe=methcol,sharp corners,
  title={\textsc{Méthode}\ifx\relax#1\relax\relax\else\textmd{ --- \itshape #1}\fi}}
\newtcolorbox{exemplebox}[1][]{boxbase,colback=excol!5,colframe=excol,
  title={\textsc{Exemple}\ifx\relax#1\relax\relax\else\textmd{ : \itshape #1}\fi}}
\newtcolorbox{piege}[1][]{boxbase,colback=attncol!6,colframe=attncol,
  title={\textsc{Piège}\ifx\relax#1\relax\relax\else\textmd{ : \itshape #1}\fi}}
\newtcolorbox{remarque}[1][]{boxbase,colback=black!4,colframe=black!45,
  title={\textsc{Remarque}\ifx\relax#1\relax\relax\else\textmd{ : \itshape #1}\fi}}
% --- boîtes EXERCICE / CORRIGÉ (noms EXACTS, auto-comptées) ---
\newtcolorbox[auto counter]{exo}[1][]{boxbase,colback=primary!4,colframe=primary,
  title={\textsc{Exercice}~\thetcbcounter\ifx\relax#1\relax\relax\else\textmd{ --- \itshape #1}\fi}}
\newtcolorbox[auto counter]{corrige}[1][]{boxbase,colback=excol!4,colframe=excol,
  title={\textsc{Corrigé}~\thetcbcounter\ifx\relax#1\relax\relax\else\textmd{ --- \itshape #1}\fi}}
\newcommand{\vv}[1]{\vec{#1}}\newcommand{\ds}{\displaystyle}
\newcommand{\R}{\mathbb{R}}\newcommand{\N}{\mathbb{N}}\newcommand{\Z}{\mathbb{Z}}
% (un \usepackage{fancyhdr}+en-tête est possible ; il est ignoré côté web)
% =========================================================================
\begin{document}
\begin{corrige}[force de Lorentz avec un angle]
\[ F=B\,|q|\,v\,\sin\alpha=0{,}4\times1{,}6\cdot10^{-19}\times2\cdot10^{6}\times\sin\ang{30}
    =0{,}4\times1{,}6\cdot10^{-19}\times2\cdot10^{6}\times0{,}5=\qty{6.4e-14}{\newton}. \]
\end{corrige}

\begin{corrige}[rayon et doublement de la vitesse]
\begin{enumerate}[label=\alph*)]
  \item $R=\dfrac{m\,v}{|q|\,B}=\dfrac{1{,}67\cdot10^{-27}\times2\cdot10^{6}}{1{,}6\cdot10^{-19}\times0{,}4}
        =\dfrac{3{,}34\cdot10^{-21}}{6{,}4\cdot10^{-20}}\approx\qty{5.2e-2}{\metre}\approx\qty{5.2}{\centi\metre}.$
  \item $R\propto v$ à $B$ constant : doubler $v$ \textbf{double} le rayon $\Rightarrow R'\approx\qty{10.4}{\centi\metre}$.
\end{enumerate}
\end{corrige}

\begin{corrige}[nature de la trajectoire d'un électron]
\begin{enumerate}[label=\alph*)]
  \item Non : la force de Lorentz est \textbf{toujours perpendiculaire} à $\vv{v}$, donc son travail est
        \textbf{nul}. Elle ne modifie pas la norme de la vitesse (mouvement \emph{uniforme}), seulement
        sa direction.
  \item Comme $\vv{v}\perp\vv{B}$ et que la vitesse reste constante, la trajectoire est un
        \textbf{cercle} parcouru à vitesse constante : un \textbf{mouvement circulaire uniforme}, dans
        le plan perpendiculaire à $\vv{B}$ (ici horizontal). (Ce n'est \emph{pas} une parabole : le poids
        est négligé.)
\end{enumerate}
\end{corrige}

\begin{corrige}[le sens pour une charge négative]
On applique d'abord la main droite comme pour une charge positive : pouce $=\vv{v}$ (droite),
index $=\vv{B}$ (haut) $\Rightarrow$ majeur \textbf{vers toi} ($\odot$). Mais la charge est
\textbf{négative} : on \textbf{inverse}. La force $\vv{F}$ \textbf{entre} donc dans la feuille,
$\otimes$.
\end{corrige}

\begin{corrige}[la force qui ne travaille pas]
\begin{enumerate}[label=\alph*)]
  \item Non : $\vv{F}\perp\vv{v}$, son travail est nul, donc l'\textbf{énergie cinétique} et la
        \textbf{norme} de la vitesse restent \textbf{constantes}.
  \item $R=\dfrac{mv}{|q|B}$ : doubler $B$ \textbf{divise le rayon par deux}. L'énergie cinétique
        $\tfrac12 mv^2$, elle, ne change \textbf{pas} (la vitesse est inchangée) : le cercle est
        simplement deux fois plus petit.
\end{enumerate}
\end{corrige}

\end{document}
